% Original assessment and worked key governed by the supplied assessment brief.
% Build from this directory; all dependencies remain in this academic-year tree.
\makeatletter
\def\input@path{{../../../preamble/}}
\makeatother
\input{preamble_exams.tex}
\input{highlights/highlights.tex}
% \GTenHighlightsfalse % Uncomment to disable instructional highlights.
\input{highlights/specific_highlights.tex}
\CoordinatedTermTwoPalette
\tikzset{gten marker/.append style={inner xsep=.6mm,inner ysep=.3mm}}
\usepackage{varwidth}
\usepackage{needspace}
% Natural cell widths, bounded locally so long descriptions can wrap.
\newcommand{\CreditDescription}[1]{\begin{varwidth}[t]{.72\linewidth}\raggedright#1\end{varwidth}}
% Match the reference practice graph treatments; the toggle also controls ink.
\colorlet{keyMassColor}{blue!75!black}
\colorlet{keyCDFColor}{cyan!65!blue!80!black}
\colorlet{keyMassFill}{GTenHLProbability!32!white}
\ifGTenHighlights\else
  \colorlet{keyMassColor}{black}
  \colorlet{keyCDFColor}{black}
  \colorlet{keyMassFill}{black!15}
\fi
\linespread{0.98}
\setlength{\parskip}{2pt}
\renewcommand{\arraystretch}{1.0}
\setlength{\emergencystretch}{2em}
\pgfplotsset{keyaxis/.style={axis lines=left,scaled ticks=false,
 xlabel={$\GTenHLValue{x}$ (\GTenHLValue{USD})},
 xticklabel={ $\GTenHLValue{\pgfmathprintnumber{\tick}}$},grid=major,grid style={gray!15},
 tick label style={font=\small},label style={font=\small},clip=false}}
\title{Global Economics Exam Solutions\\Grade 10 --- Term 2 --- C1 and C2}
\author{}
\date{}
\makeatletter
\renewcommand{\@maketitle}{\begin{center}{\Large\bfseries\@title\par}\end{center}}
\renewcommand\subsection{\@startsection{subsection}{2}{\z@}{-1.4ex plus -.3ex}{.7ex}{\normalfont\large\bfseries}}
\makeatother
\hypersetup{pdftitle={Grade 10 Term 2 C1 and C2 Exam Solutions},pdfauthor={Global Economics}}
\begin{document}
\maketitle
\small
\setlength{\topsep}{3pt}
\setcounter{tocdepth}{1}
\phantomsection
\label{toc}
\tableofcontents
\section*{Learning objective}
Model discrete monetary outcomes in economic and financial contexts; derive and verify probability mass functions (PMFs), and construct and interpret cumulative distribution functions (CDFs) with tables and graphs.
\section*{Criteria assessment}
This standalone key authors the original questions required by the governing assessment brief and reproduces each task before its worked solution. Questions 1--2 assess C1; Questions 3--6 assess C2. Show reasoning and exact probabilities. All contextual rules are stated in ordinary language.

\textbf{Independent passing thresholds: C1 requires at least 4 out of 6 points; C2 requires at least 4 out of 14 points. Points cannot transfer between criteria.} The overall available total is 20 points. An overall total alone does not establish a pass in both criteria. These are the passing thresholds for this assessment; no cumulative policy is specified.
\section*{Evaluated criteria}
\begin{description}
\item[C1] Characterises discrete random variables in economic and financial contexts by identifying their support, probability mass function and possible outcomes.
\item[C2] Interprets PMFs and CDFs by verifying their defining properties and using them to calculate point, interval, cumulative, and complementary probabilities in financial contexts.
\end{description}

\begin{center}
\begin{tabular}{|c|c|c|}
	\hline\textbf{\textnumero} & \textbf{C1} & \textbf{C2} \\ \hline
	\ref*{sec:q1} & 3 & --- \\\hline
	\ref*{sec:q2} & 3 & --- \\\hline
	\ref*{sec:q3} & --- & 2 \\\hline
	\ref*{sec:q4} & --- & 2 \\\hline
	\ref*{sec:q5} & --- & 4 \\\hline
	\ref*{sec:q6} & --- & 6 \\\hline
	\textbf{Total} & \textbf{6} & \textbf{14} \\\hline
\end{tabular}
\end{center}

\textbf{Marking policy.} Question and criterion totals are fixed by the brief. Every task-level split and partial-credit breakdown below is proposed. Award stated fractions for demonstrated components, including valid alternative methods. Allow consistent follow-through after one arithmetic error when subsequent reasoning is correct. In C1 Task 2 (Questions 1--2) and C2 Task 1 (Questions 3--6), positivity (nonnegativity) and normalization checks carry no separate points; omitting these checks does not reduce the task score. Solving the normalization equation to derive probabilities in Question 5 remains assessed, as do its relationship checks. CDF verification retains its stated credit and must be demonstrated. Cap each task, question and criterion at its allocation.


\clearpage
\section*{Highlight legend}
\phantomsection\label{tab:highlight-legend}
\begin{center}
\begin{tabular}{@{}lll@{}}\toprule
\textbf{Symbol or notation} & \textbf{Concept} & \textbf{Brief meaning}\\\midrule
$\GTenHLOutcome{\omega}$ & \GTenHLOutcome{Outcome} & One complete result or category.\\
$\GTenHLValue{X},\GTenHLValue{x},\GTenHLValue{S_X}$ & \GTenHLValue{Variable and support} & Values and cutoffs in \GTenHLValue{USD}.\\
$\GTenHLSecond{N},\GTenHLSecond{Y}$ & \GTenHLSecond{Other quantities} & Counts; adjusted payout in USD.\\
$\GTenHLProbability{P},\GTenHLProbability{p_X(x)}$ & \GTenHLProbability{Probability / PMF} & Unitless likelihood or point mass.\\
$\HLCumulativeProbability{F_X(x)}$ & \HLCumulativeProbability{Cumulative probability / CDF} & Mass at or below a cutoff.\\
$\GTenHLParameter{a}$ & \GTenHLParameter{Unknown parameter} & Baseline mass to determine in Question 5.\\
\bottomrule\end{tabular}
\end{center}

\sectionScored{Online order contribution}{C1-\sml{3}}{sec:q1}

A small seller handles three orders, in first, second and third order. Each order independently has exactly two mutually exclusive possible outcomes: the \GTenHLOutcome{buyer keeps it} with probability \GTenHLProbability{0.6}, or \GTenHLOutcome{cancels it} with probability \GTenHLProbability{0.4}. These probabilities apply to every order. Each kept order contributes \GTenHLValue{5 USD} and each cancellation contributes no money. A fixed handling cost of \GTenHLValue{4 USD} is then deducted once from the combined contribution. There are no other receipts or costs. Record the \GTenHLValue{net contribution} for the three orders.

\Needspace{8\baselineskip}
\subsection*{Task 1 --- question and worked solution (proposed allocation) \score{1}}
\label{sol:q1t1}

\noindent\textbf{\boldmath List every complete outcome in order and calculate its probability.}\par\nopagebreak

Write $\GTenHLOutcome{K}$ for \GTenHLOutcome{kept} and $\GTenHLOutcome{C}$ for \GTenHLOutcome{cancelled}. There are $2^{3}=8$ ordered complete outcomes. Independence permits multiplication within each triple; it does not make the triples equally likely.

\begin{center}
\begin{tabular}{ccc}\toprule
Ordered outcome & Product & Exact probability\\\midrule
$\GTenHLOutcome{\mathrm{KKK}}$ & $\GTenHLProbability{0.6}\cdot\GTenHLProbability{0.6}\cdot\GTenHLProbability{0.6}$ & $\GTenHLProbability{27/125}=\GTenHLProbability{0.216}$\\
$\GTenHLOutcome{\mathrm{KKC}}$ & $\GTenHLProbability{0.6}\cdot\GTenHLProbability{0.6}\cdot\GTenHLProbability{0.4}$ & $\GTenHLProbability{18/125}=\GTenHLProbability{0.144}$\\
$\GTenHLOutcome{\mathrm{KCK}}$ & $\GTenHLProbability{0.6}\cdot\GTenHLProbability{0.4}\cdot\GTenHLProbability{0.6}$ & $\GTenHLProbability{18/125}=\GTenHLProbability{0.144}$\\
$\GTenHLOutcome{\mathrm{KCC}}$ & $\GTenHLProbability{0.6}\cdot\GTenHLProbability{0.4}\cdot\GTenHLProbability{0.4}$ & $\GTenHLProbability{12/125}=\GTenHLProbability{0.096}$\\
$\GTenHLOutcome{\mathrm{CKK}}$ & $\GTenHLProbability{0.4}\cdot\GTenHLProbability{0.6}\cdot\GTenHLProbability{0.6}$ & $\GTenHLProbability{18/125}=\GTenHLProbability{0.144}$\\
$\GTenHLOutcome{\mathrm{CKC}}$ & $\GTenHLProbability{0.4}\cdot\GTenHLProbability{0.6}\cdot\GTenHLProbability{0.4}$ & $\GTenHLProbability{12/125}=\GTenHLProbability{0.096}$\\
$\GTenHLOutcome{\mathrm{CCK}}$ & $\GTenHLProbability{0.4}\cdot\GTenHLProbability{0.4}\cdot\GTenHLProbability{0.6}$ & $\GTenHLProbability{12/125}=\GTenHLProbability{0.096}$\\
$\GTenHLOutcome{\mathrm{CCC}}$ & $\GTenHLProbability{0.4}\cdot\GTenHLProbability{0.4}\cdot\GTenHLProbability{0.4}$ & $\GTenHLProbability{8/125}=\GTenHLProbability{0.064}$\\
\bottomrule\end{tabular}
\end{center}

The eight probabilities sum to $(\GTenHLProbability{3/5}+\GTenHLProbability{2/5})^{3}=\GTenHLProbability{1}$.

\par\smallskip
\noindent\begin{minipage}{\linewidth}\centering
\textbf{Proposed partial credit.}\par\smallskip
\begin{tabular}{@{}l r@{}}\toprule
\CreditDescription{\textbf{Assessed step or criterion}} & \textbf{Points}\\\midrule
\CreditDescription{All eight ordered outcomes} & 0.5\\
\CreditDescription{Their probabilities using independent products} & 0.5\\
\midrule
\CreditDescription{\textbf{Total}} & \textbf{1}\\
\bottomrule\end{tabular}
\end{minipage}\par\smallskip

\par\vspace{3mm}\secInicio[][sec:q1]

\Needspace{8\baselineskip}
\subsection*{Task 2 --- question and worked solution (proposed allocation) \score{1}}
\label{sol:q1t2}

\noindent\textbf{\boldmath Define the \GTenHLValue{random variable} and its units. Map every complete outcome to a value, identify the \GTenHLValue{support}, and construct and verify the grouped \GTenHLProbability{PMF}.}\par\nopagebreak

Let $\GTenHLValue{X}$ be the \GTenHLValue{net contribution} from the three orders, in \GTenHLValue{USD}. If $\GTenHLSecond{N}$ is the number kept, $\GTenHLValue{X}=\GTenHLValue{5}\GTenHLSecond{N}-\GTenHLValue{4}$. The table maps all eight outcomes and groups those with the same \GTenHLValue{net contribution}.

\begin{center}
\begin{tabular}{ccc}\toprule
$\GTenHLValue{x}$ (\GTenHLValue{USD}) & Complete outcomes & Grouped probability\\\midrule
$-\GTenHLValue{4}$ & $\GTenHLOutcome{\mathrm{CCC}}$ & $\GTenHLProbability{8/125}=\GTenHLProbability{8/125}$\\
$\GTenHLValue{1}$ & $\GTenHLOutcome{\mathrm{KCC}},\GTenHLOutcome{\mathrm{CKC}},\GTenHLOutcome{\mathrm{CCK}}$ & $\GTenHLProbability{12/125}+\GTenHLProbability{12/125}+\GTenHLProbability{12/125}=\GTenHLProbability{36/125}$\\
$\GTenHLValue{6}$ & $\GTenHLOutcome{\mathrm{KKC}},\GTenHLOutcome{\mathrm{KCK}},\GTenHLOutcome{\mathrm{CKK}}$ & $\GTenHLProbability{18/125}+\GTenHLProbability{18/125}+\GTenHLProbability{18/125}=\GTenHLProbability{54/125}$\\
$\GTenHLValue{11}$ & $\GTenHLOutcome{\mathrm{KKK}}$ & $\GTenHLProbability{27/125}=\GTenHLProbability{27/125}$\\
\bottomrule\end{tabular}
\end{center}

Thus $\GTenHLValue{S_X}=\{-\GTenHLValue{4},\GTenHLValue{1},\GTenHLValue{6},\GTenHLValue{11}\}$ \GTenHLValue{USD}. The \GTenHLProbability{PMF} is the last column.
\Needspace{5\baselineskip}\textbf{PMF checks.}
\begin{itemize}
\item Nonnegativity: all listed masses are positive.
\item Normalization: $(\GTenHLProbability{8}+\GTenHLProbability{36}+\GTenHLProbability{54}+\GTenHLProbability{27})/\GTenHLProbability{125}=\GTenHLProbability{1}$.
\item Outside the \GTenHLValue{support}: the \GTenHLProbability{PMF} is zero elsewhere.
\end{itemize}

A negative value means the handling cost exceeds receipts.

\textbf{Scoring note.} These explanatory positivity and normalization checks carry no separate points; omitting them does not reduce the task score.

\par\smallskip
\noindent\begin{minipage}{\linewidth}\centering
\textbf{Proposed partial credit.}\par\smallskip
\begin{tabular}{@{}l r@{}}\toprule
\CreditDescription{\textbf{Assessed step or criterion}} & \textbf{Points}\\\midrule
\CreditDescription{The variable and units} & 0.125\\
\CreditDescription{All outcome mappings} & 0.25\\
\CreditDescription{Support} & 0.125\\
\CreditDescription{Grouped PMF masses} & 0.5\\
\midrule
\CreditDescription{\textbf{Total}} & \textbf{1}\\
\bottomrule\end{tabular}
\end{minipage}\par\smallskip

\par\vspace{3mm}\secInicio[][sec:q1]

\Needspace{8\baselineskip}
\subsection*{Task 3 --- question and worked solution (proposed allocation) \score{1}}
\label{sol:q1t3}

\noindent\textbf{\boldmath Calculate and interpret the probability that the \GTenHLValue{net contribution} is exactly \GTenHLValue{6 USD}.}\par\nopagebreak

From the grouped \GTenHLProbability{PMF}, $P(\GTenHLValue{X}=\GTenHLValue{6})=\GTenHLProbability{p_X(6)}=\GTenHLProbability{54/125}=\GTenHLProbability{0.432}=\GTenHLProbability{43.2\%}$. The seller earns a \GTenHLValue{net contribution} of exactly \GTenHLValue{6 USD} with probability \GTenHLProbability{43.2\%}; this happens exactly when two orders are kept and one is cancelled.

\par\smallskip
\noindent\begin{minipage}{\linewidth}\centering
\textbf{Proposed partial credit.}\par\smallskip
\begin{tabular}{@{}l r@{}}\toprule
\CreditDescription{\textbf{Assessed step or criterion}} & \textbf{Points}\\\midrule
\CreditDescription{Selecting the correct mass at $X=6$} & 0.5\\
\CreditDescription{The contextual interpretation} & 0.5\\
\midrule
\CreditDescription{\textbf{Total}} & \textbf{1}\\
\bottomrule\end{tabular}
\end{minipage}\par\smallskip

Question total: 3 points for C1.

\par\vspace{3mm}\secInicio[][sec:q1]

\clearpage
\sectionScored{Two supplier loyalty credits}{C1-\sml{3}}{sec:q2}

A retailer receives one loyalty credit from each of two suppliers. The suppliers award credits independently. Each has exactly three mutually exclusive possible outcomes: no credit, a \GTenHLValue{3 USD} credit, or an \GTenHLValue{8 USD} credit. For the first supplier these probabilities are \GTenHLProbability{0.2}, \GTenHLProbability{0.5} and \GTenHLProbability{0.3}, respectively. For the second supplier they are \GTenHLProbability{0.4}, \GTenHLProbability{0.4} and \GTenHLProbability{0.2}, respectively. Add both credits; there are no other adjustments. Record the total credit, keeping first-supplier then second-supplier order.

\Needspace{8\baselineskip}
\subsection*{Task 1 --- question and worked solution (proposed allocation) \score{1}}
\label{sol:q2t1}

\noindent\textbf{\boldmath List every complete ordered outcome and calculate its probability.}\par\nopagebreak

Each supplier has three possible outcomes, giving $3^{2}=9$ ordered pairs. Write a pair as its two credit amounts, in \GTenHLValue{USD}. Multiply the appropriate probabilities by independence.

\begin{center}
\begin{tabular}{ccc}\toprule
Ordered credit outcome (\GTenHLValue{USD}) & Product & Exact probability\\\midrule
$(\GTenHLOutcome{0},\GTenHLOutcome{0})$ & $\GTenHLProbability{0.2}\cdot\GTenHLProbability{0.4}$ & $\GTenHLProbability{2/25}=\GTenHLProbability{0.08}$\\
$(\GTenHLOutcome{0},\GTenHLOutcome{3})$ & $\GTenHLProbability{0.2}\cdot\GTenHLProbability{0.4}$ & $\GTenHLProbability{2/25}=\GTenHLProbability{0.08}$\\
$(\GTenHLOutcome{0},\GTenHLOutcome{8})$ & $\GTenHLProbability{0.2}\cdot\GTenHLProbability{0.2}$ & $\GTenHLProbability{1/25}=\GTenHLProbability{0.04}$\\
$(\GTenHLOutcome{3},\GTenHLOutcome{0})$ & $\GTenHLProbability{0.5}\cdot\GTenHLProbability{0.4}$ & $\GTenHLProbability{1/5}=\GTenHLProbability{0.20}$\\
$(\GTenHLOutcome{3},\GTenHLOutcome{3})$ & $\GTenHLProbability{0.5}\cdot\GTenHLProbability{0.4}$ & $\GTenHLProbability{1/5}=\GTenHLProbability{0.20}$\\
$(\GTenHLOutcome{3},\GTenHLOutcome{8})$ & $\GTenHLProbability{0.5}\cdot\GTenHLProbability{0.2}$ & $\GTenHLProbability{1/10}=\GTenHLProbability{0.10}$\\
$(\GTenHLOutcome{8},\GTenHLOutcome{0})$ & $\GTenHLProbability{0.3}\cdot\GTenHLProbability{0.4}$ & $\GTenHLProbability{3/25}=\GTenHLProbability{0.12}$\\
$(\GTenHLOutcome{8},\GTenHLOutcome{3})$ & $\GTenHLProbability{0.3}\cdot\GTenHLProbability{0.4}$ & $\GTenHLProbability{3/25}=\GTenHLProbability{0.12}$\\
$(\GTenHLOutcome{8},\GTenHLOutcome{8})$ & $\GTenHLProbability{0.3}\cdot\GTenHLProbability{0.2}$ & $\GTenHLProbability{3/50}=\GTenHLProbability{0.06}$\\
\bottomrule\end{tabular}
\end{center}

Exhaustiveness check: $(\GTenHLProbability{0.2}+\GTenHLProbability{0.5}+\GTenHLProbability{0.3})(\GTenHLProbability{0.4}+\GTenHLProbability{0.4}+\GTenHLProbability{0.2})=\GTenHLProbability{1}$.

\par\smallskip
\noindent\begin{minipage}{\linewidth}\centering
\textbf{Proposed partial credit.}\par\smallskip
\begin{tabular}{@{}l r@{}}\toprule
\CreditDescription{\textbf{Assessed step or criterion}} & \textbf{Points}\\\midrule
\CreditDescription{All nine ordered pairs} & 0.5\\
\CreditDescription{Independent product probabilities} & 0.5\\
\midrule
\CreditDescription{\textbf{Total}} & \textbf{1}\\
\bottomrule\end{tabular}
\end{minipage}\par\smallskip

\par\vspace{3mm}\secInicio[][sec:q2]

\Needspace{8\baselineskip}
\subsection*{Task 2 --- question and worked solution (proposed allocation) \score{1}}
\label{sol:q2t2}

\noindent\textbf{\boldmath Define the \GTenHLValue{random variable} and units; map every outcome to a value; identify the \GTenHLValue{support}; construct and verify the grouped \GTenHLProbability{PMF}.}\par\nopagebreak

Let $\GTenHLValue{X}$ be the \GTenHLValue{total loyalty credit} from the two suppliers, in \GTenHLValue{USD}. Add the two amounts in each pair. The complete mapping and grouped \GTenHLProbability{PMF} are:

\begin{center}
\begin{tabular}{ccc}\toprule
$\GTenHLValue{x}$ (\GTenHLValue{USD}) & Ordered outcomes & Grouped probability\\\midrule
$\GTenHLValue{0}$ & $(\GTenHLOutcome{0},\GTenHLOutcome{0})$ & $\GTenHLProbability{2/25}=\GTenHLProbability{2/25}$\\
$\GTenHLValue{3}$ & $(\GTenHLOutcome{0},\GTenHLOutcome{3}),(\GTenHLOutcome{3},\GTenHLOutcome{0})$ & $\GTenHLProbability{2/25}+\GTenHLProbability{1/5}=\GTenHLProbability{7/25}$\\
$\GTenHLValue{6}$ & $(\GTenHLOutcome{3},\GTenHLOutcome{3})$ & $\GTenHLProbability{1/5}=\GTenHLProbability{1/5}$\\
$\GTenHLValue{8}$ & $(\GTenHLOutcome{0},\GTenHLOutcome{8}),(\GTenHLOutcome{8},\GTenHLOutcome{0})$ & $\GTenHLProbability{1/25}+\GTenHLProbability{3/25}=\GTenHLProbability{4/25}$\\
$\GTenHLValue{11}$ & $(\GTenHLOutcome{3},\GTenHLOutcome{8}),(\GTenHLOutcome{8},\GTenHLOutcome{3})$ & $\GTenHLProbability{1/10}+\GTenHLProbability{3/25}=\GTenHLProbability{11/50}$\\
$\GTenHLValue{16}$ & $(\GTenHLOutcome{8},\GTenHLOutcome{8})$ & $\GTenHLProbability{3/50}=\GTenHLProbability{3/50}$\\
\bottomrule\end{tabular}
\end{center}

The \GTenHLValue{support} is $\GTenHLValue{S_X}=\{\GTenHLValue{0},\GTenHLValue{3},\GTenHLValue{6},\GTenHLValue{8},\GTenHLValue{11},\GTenHLValue{16}\}$ \GTenHLValue{USD}, with the grouped masses shown above.
\Needspace{5\baselineskip}\textbf{PMF checks.}
\begin{itemize}
\item Nonnegativity: all listed masses are positive.
\item Normalization: $(\GTenHLProbability{4}+\GTenHLProbability{14}+\GTenHLProbability{10}+\GTenHLProbability{8}+\GTenHLProbability{11}+\GTenHLProbability{3})/\GTenHLProbability{50}=\GTenHLProbability{50/50}=\GTenHLProbability{1}$.
\item Outside the \GTenHLValue{support}: the \GTenHLProbability{PMF} has zero mass elsewhere.
\end{itemize}

Reversing a pair can preserve its monetary value while changing its probability.

\textbf{Scoring note.} These explanatory positivity and normalization checks carry no separate points; omitting them does not reduce the task score.

\par\smallskip
\noindent\begin{minipage}{\linewidth}\centering
\textbf{Proposed partial credit.}\par\smallskip
\begin{tabular}{@{}l r@{}}\toprule
\CreditDescription{\textbf{Assessed step or criterion}} & \textbf{Points}\\\midrule
\CreditDescription{Variable and units} & 0.125\\
\CreditDescription{All mappings} & 0.25\\
\CreditDescription{Support} & 0.125\\
\CreditDescription{Grouped masses} & 0.5\\
\midrule
\CreditDescription{\textbf{Total}} & \textbf{1}\\
\bottomrule\end{tabular}
\end{minipage}\par\smallskip

\par\vspace{3mm}\secInicio[][sec:q2]

\Needspace{8\baselineskip}
\subsection*{Task 3 --- question and worked solution (proposed allocation) \score{1}}
\label{sol:q2t3}

\noindent\textbf{\boldmath Calculate and interpret the probability that the total credit is exactly \GTenHLValue{11 USD}.}\par\nopagebreak

$P(\GTenHLValue{X}=\GTenHLValue{11})=\GTenHLProbability{p_X(11)}=\GTenHLProbability{11/50}=\GTenHLProbability{0.22}=\GTenHLProbability{22\%}$. The retailer receives a total credit of exactly \GTenHLValue{11 USD} with probability \GTenHLProbability{22\%}; the qualifying pairs are $(\GTenHLValue{3},\GTenHLValue{8})$ and $(\GTenHLValue{8},\GTenHLValue{3})$.

\par\smallskip
\noindent\begin{minipage}{\linewidth}\centering
\textbf{Proposed partial credit.}\par\smallskip
\begin{tabular}{@{}l r@{}}\toprule
\CreditDescription{\textbf{Assessed step or criterion}} & \textbf{Points}\\\midrule
\CreditDescription{Selecting the correct mass at $X=11$} & 0.5\\
\CreditDescription{The contextual interpretation} & 0.5\\
\midrule
\CreditDescription{\textbf{Total}} & \textbf{1}\\
\bottomrule\end{tabular}
\end{minipage}\par\smallskip

Question total: 3 points for C1. Criterion C1 total: 6 points; independent pass: at least 4 points.

\par\vspace{3mm}\secInicio[][sec:q2]

\clearpage
\sectionScored{Local parcel delivery fees}{C2-\sml{2}}{sec:q3}

\textbf{Basic Type A.} A cooperative sends one parcel using exactly one of eight exhaustive, mutually exclusive delivery methods: \GTenHLOutcome{walking courier}, \GTenHLOutcome{bicycle}, \GTenHLOutcome{cargo bicycle}, \GTenHLOutcome{motorcycle}, car, van, \GTenHLOutcome{refrigerated van} or truck. Their probabilities are \GTenHLProbability{0.10}, \GTenHLProbability{0.10}, \GTenHLProbability{0.15}, \GTenHLProbability{0.20}, \GTenHLProbability{0.15}, \GTenHLProbability{0.10}, \GTenHLProbability{0.10} and \GTenHLProbability{0.10}, respectively. The corresponding fixed handling fees are \GTenHLValue{0}, \GTenHLValue{0}, \GTenHLValue{2}, \GTenHLValue{4}, \GTenHLValue{7}, \GTenHLValue{11}, \GTenHLValue{16} and \GTenHLValue{23 USD}, respectively. No other handling charges apply. Record the \GTenHLValue{handling fee} for the parcel.

\Needspace{8\baselineskip}
\subsection*{Task 1 --- question and worked solution (proposed allocation) \score{0.5}}
\label{sol:q3t1}

\noindent\textbf{\boldmath Define the \GTenHLValue{random variable} and \GTenHLValue{support}; derive and verify the grouped \GTenHLProbability{PMF}; produce its table and a labeled graph with actual monetary spacing.}\par\nopagebreak

Let $\GTenHLValue{X}$ be the \GTenHLValue{handling fee} for one parcel, in \GTenHLValue{USD}. Map the supplied categories to payments:

\begin{center}
\begin{tabular}{ccc}\toprule
Category & Given probability & $\GTenHLValue{X}$ (\GTenHLValue{USD})\\\midrule
\GTenHLOutcome{Walking courier} & $\GTenHLProbability{0.10}=\GTenHLProbability{2/20}$ & $\GTenHLValue{0}$\\
\GTenHLOutcome{Bicycle} & $\GTenHLProbability{0.10}=\GTenHLProbability{2/20}$ & $\GTenHLValue{0}$\\
\GTenHLOutcome{Cargo bicycle} & $\GTenHLProbability{0.15}=\GTenHLProbability{3/20}$ & $\GTenHLValue{2}$\\
\GTenHLOutcome{Motorcycle} & $\GTenHLProbability{0.20}=\GTenHLProbability{4/20}$ & $\GTenHLValue{4}$\\
\GTenHLOutcome{Car} & $\GTenHLProbability{0.15}=\GTenHLProbability{3/20}$ & $\GTenHLValue{7}$\\
\GTenHLOutcome{Van} & $\GTenHLProbability{0.10}=\GTenHLProbability{2/20}$ & $\GTenHLValue{11}$\\
\GTenHLOutcome{Refrigerated van} & $\GTenHLProbability{0.10}=\GTenHLProbability{2/20}$ & $\GTenHLValue{16}$\\
\GTenHLOutcome{Truck} & $\GTenHLProbability{0.10}=\GTenHLProbability{2/20}$ & $\GTenHLValue{23}$\\
\bottomrule\end{tabular}
\end{center}

The walking courier and bicycle are distinct outcomes sharing a zero fee, so $\GTenHLProbability{p_X(0)}=\GTenHLProbability{0.10}+\GTenHLProbability{0.10}=\GTenHLProbability{0.20}=\GTenHLProbability{4/20}$. Every other method has its own fee and retains its supplied probability.

$\GTenHLValue{S_X}=\{\GTenHLValue{0},\GTenHLValue{2},\GTenHLValue{4},\GTenHLValue{7},\GTenHLValue{11},\GTenHLValue{16},\GTenHLValue{23}\}$ \GTenHLValue{USD}. There are eight elementary outcomes but seven distinct monetary values. Add probabilities of mutually exclusive categories sharing a value, because the events cannot occur together.

\begin{center}
\begin{minipage}{.26\linewidth}\centering
\begin{tabular}{cc}\toprule
$\GTenHLValue{x}$ (\GTenHLValue{USD}) & $\GTenHLProbability{p_X(x)}$\\\midrule
$\GTenHLValue{0}$ & $\GTenHLProbability{4/20}$\\
$\GTenHLValue{2}$ & $\GTenHLProbability{3/20}$\\
$\GTenHLValue{4}$ & $\GTenHLProbability{4/20}$\\
$\GTenHLValue{7}$ & $\GTenHLProbability{3/20}$\\
$\GTenHLValue{11}$ & $\GTenHLProbability{2/20}$\\
$\GTenHLValue{16}$ & $\GTenHLProbability{2/20}$\\
$\GTenHLValue{23}$ & $\GTenHLProbability{2/20}$\\
\bottomrule\end{tabular}
\end{minipage}\hfill
\begin{minipage}{.70\linewidth}\centering
\begin{tikzpicture}
\begin{axis}[keyaxis,width=.94\linewidth,height=5cm,xmin=-2,xmax=25,ymin=0,ymax=.35,xtick={0,2,4,7,11,16,23},xticklabel style={rotate=45,anchor=north east},ylabel={$\GTenHLProbability{p_X(x)}$},yticklabel={$\GTenHLProbability{\pgfmathprintnumber{\tick}}$},ytick={0,.05,.1,.15,.2,.25,.3}]
\draw[fill=keyMassFill,draw=keyMassColor] (axis cs:-0.45,0) rectangle (axis cs:0.45,4/20);
\draw[fill=keyMassFill,draw=keyMassColor] (axis cs:1.55,0) rectangle (axis cs:2.45,3/20);
\draw[fill=keyMassFill,draw=keyMassColor] (axis cs:3.55,0) rectangle (axis cs:4.45,4/20);
\draw[fill=keyMassFill,draw=keyMassColor] (axis cs:6.55,0) rectangle (axis cs:7.45,3/20);
\draw[fill=keyMassFill,draw=keyMassColor] (axis cs:10.55,0) rectangle (axis cs:11.45,2/20);
\draw[fill=keyMassFill,draw=keyMassColor] (axis cs:15.55,0) rectangle (axis cs:16.45,2/20);
\draw[fill=keyMassFill,draw=keyMassColor] (axis cs:22.55,0) rectangle (axis cs:23.45,2/20);
\end{axis}\end{tikzpicture}
\end{minipage}
\end{center}

\Needspace{5\baselineskip}\textbf{PMF checks.}
\begin{itemize}
\item Nonnegativity: every listed mass is strictly positive.
\item Normalization: $\GTenHLProbability{4}+\GTenHLProbability{3}+\GTenHLProbability{4}+\GTenHLProbability{3}+\GTenHLProbability{2}+\GTenHLProbability{2}+\GTenHLProbability{2}=\GTenHLProbability{20}$, so the masses sum to $\GTenHLProbability{20/20}=\GTenHLProbability{1}$.
\item Outside the \GTenHLValue{support}: set $\GTenHLProbability{p_X(x)}=\GTenHLProbability{0}$ there. The graph places bars at the actual monetary values; gaps carry no mass.
\end{itemize}

\textbf{Scoring note.} These explanatory positivity and normalization checks carry no separate points; omitting them does not reduce the task score.

\par\smallskip
\noindent\begin{minipage}{\linewidth}\centering
\textbf{Proposed partial credit.}\par\smallskip
\begin{tabular}{@{}l r@{}}\toprule
\CreditDescription{\textbf{Assessed step or criterion}} & \textbf{Points}\\\midrule
\CreditDescription{Variable, units and support} & 0.125\\
\CreditDescription{Category mapping, grouping and the PMF table} & 0.25\\
\CreditDescription{The labeled, correctly spaced graph} & 0.125\\
\midrule
\CreditDescription{\textbf{Total}} & \textbf{0.5}\\
\bottomrule\end{tabular}
\end{minipage}\par\smallskip

\par\vspace{3mm}\secInicio[][sec:q3]

\clearpage
\textbf{Question 3 --- continued}\quad\hyperref[sec:q3]{Original description}

\Needspace{8\baselineskip}
\subsection*{Task 2 --- question and worked solution (proposed allocation) \score{1}}
\label{sol:q3t2}

\noindent\textbf{\boldmath Construct the complete \HLCumulativeProbability{CDF} table, including both tails, and draw a labeled step graph. Use open and filled endpoints correctly and preserve actual monetary spacing. Verify the defining \HLCumulativeProbability{CDF} properties and that its jumps equal the \GTenHLProbability{PMF} masses.}\par\nopagebreak

For $\HLCumulativeProbability{F_X(x)}=P(\GTenHLValue{X}\leq \GTenHLValue{x})$, add masses in increasing fee order. The successive cumulative numerators are $\HLCumulativeProbability{4},\HLCumulativeProbability{7},\HLCumulativeProbability{11},\HLCumulativeProbability{14},\HLCumulativeProbability{16},\HLCumulativeProbability{18},\HLCumulativeProbability{20}$, all over $\HLCumulativeProbability{20}$.

\begin{center}
\begin{tabular}{cc}\toprule
Cutoff interval (\GTenHLValue{USD}) & $\HLCumulativeProbability{F_X(x)}$\\\midrule
$\GTenHLValue{x}<\GTenHLValue{0}$ & $\HLCumulativeProbability{0}$\\
$\GTenHLValue{0}\leq \GTenHLValue{x}<\GTenHLValue{2}$ & $\HLCumulativeProbability{4/20}$\\
$\GTenHLValue{2}\leq \GTenHLValue{x}<\GTenHLValue{4}$ & $\HLCumulativeProbability{7/20}$\\
$\GTenHLValue{4}\leq \GTenHLValue{x}<\GTenHLValue{7}$ & $\HLCumulativeProbability{11/20}$\\
$\GTenHLValue{7}\leq \GTenHLValue{x}<\GTenHLValue{11}$ & $\HLCumulativeProbability{14/20}$\\
$\GTenHLValue{11}\leq \GTenHLValue{x}<\GTenHLValue{16}$ & $\HLCumulativeProbability{16/20}$\\
$\GTenHLValue{16}\leq \GTenHLValue{x}<\GTenHLValue{23}$ & $\HLCumulativeProbability{18/20}$\\
$\GTenHLValue{x}\geq \GTenHLValue{23}$ & $\HLCumulativeProbability{1}$\\
\bottomrule\end{tabular}
\end{center}

\begin{center}\begin{tikzpicture}
\begin{axis}[keyaxis,width=.86\linewidth,height=7cm,xmin=-2,xmax=25,ymin=-.07,ymax=1.1,xtick={0,2,4,7,11,16,23},xticklabel style={rotate=45,anchor=north east},ytick={0,.25,.5,.75,1},ylabel={$\HLCumulativeProbability{F_X(x)}$},yticklabel={$\HLCumulativeProbability{\pgfmathprintnumber{\tick}}$}]
\addplot[keyCDFColor,thick,<-] coordinates {(-2,0)(0,0)};
\addplot[keyCDFColor,thick] coordinates {(0,4/20)(2,4/20)};
\addplot[keyCDFColor,thick] coordinates {(2,7/20)(4,7/20)};
\addplot[keyCDFColor,thick] coordinates {(4,11/20)(7,11/20)};
\addplot[keyCDFColor,thick] coordinates {(7,14/20)(11,14/20)};
\addplot[keyCDFColor,thick] coordinates {(11,16/20)(16,16/20)};
\addplot[keyCDFColor,thick] coordinates {(16,18/20)(23,18/20)};
\addplot[keyCDFColor,thick,->] coordinates {(23,1)(25,1)};
\addplot[keyCDFColor,only marks,mark=*,mark size=1.6pt,mark options={fill=white}] coordinates {(0,0/20)(2,4/20)(4,7/20)(7,11/20)(11,14/20)(16,16/20)(23,18/20)};
\addplot[keyCDFColor,only marks,mark=*,mark size=1.6pt] coordinates {(0,4/20)(2,7/20)(4,11/20)(7,14/20)(11,16/20)(16,18/20)(23,20/20)};
\end{axis}\end{tikzpicture}\end{center}

\textbf{CDF checks.}
\begin{itemize}
\item Domain: the table covers every real cutoff.
\item Range: all tabulated heights lie in $[\HLCumulativeProbability{0},\HLCumulativeProbability{1}]$.
\item Monotonicity: the tabulated heights never decrease.
\item Right-continuity: each horizontal segment includes its left endpoint and excludes its right endpoint; the filled point gives the value at a jump.
\item Lower-tail limit: the left arrow extends to $-\infty$, where the limit is 0.
\item Upper-tail limit: the right arrow extends to $+\infty$, where the limit is 1.
\item Jump--mass agreement: in \GTenHLValue{support} order, subtracting successive cumulative numerators (starting at zero) gives $\GTenHLProbability{4},\GTenHLProbability{3},\GTenHLProbability{4},\GTenHLProbability{3},\GTenHLProbability{2},\GTenHLProbability{2},\GTenHLProbability{2}$. Hence each jump is its \GTenHLProbability{PMF} mass over $\GTenHLProbability{20}$: $\HLCumulativeProbability{F_X(a)}-\HLCumulativeProbability{F_X(a^-)}=\GTenHLProbability{p_X(a)}$. At unsupported cutoffs both the jump and the mass are zero.
\end{itemize}

\par\smallskip
\noindent\begin{minipage}{\linewidth}\centering
\textbf{Proposed partial credit.}\par\smallskip
\begin{tabular}{@{}l r@{}}\toprule
\CreditDescription{\textbf{Assessed step or criterion}} & \textbf{Points}\\\midrule
\CreditDescription{Cumulative sums} & 0.25\\
\CreditDescription{The complete interval table and tails} & 0.25\\
\CreditDescription{The labeled, correctly spaced step graph and endpoints} & 0.25\\
\CreditDescription{The CDF and jump checks} & 0.25\\
\midrule
\CreditDescription{\textbf{Total}} & \textbf{1}\\
\bottomrule\end{tabular}
\end{minipage}\par\smallskip

\par\vspace{3mm}\secInicio[][sec:q3]

\clearpage
\textbf{Question 3 --- continued}\quad\hyperref[sec:q3]{Original description}

\Needspace{8\baselineskip}
\subsection*{Task 3 (c) --- question and worked solution (proposed allocation) \score{0.5}}
\label{sol:q3t3}

\noindent\textbf{\boldmath Calculate and interpret (i) the probability that the \GTenHLValue{handling fee} is less than \GTenHLValue{7 USD} and (ii) the probability that it is at least \GTenHLValue{16 USD}.}\par\nopagebreak

\textbf{Allocation:} (i) 0.25 points; (ii) 0.25 points.

\textbf{(i) Less than \GTenHLValue{7 USD}.} Include only the \GTenHLValue{support} values $\GTenHLValue{0},\GTenHLValue{2},\GTenHLValue{4}$; the mass at \GTenHLValue{7 USD} is excluded because the inequality is strict:
\[
P(\GTenHLValue{X}<\GTenHLValue{7})=\GTenHLProbability{p_X(0)}+\GTenHLProbability{p_X(2)}+\GTenHLProbability{p_X(4)}
=\frac{\GTenHLProbability{4}+\GTenHLProbability{3}+\GTenHLProbability{4}}{\GTenHLProbability{20}}=\GTenHLProbability{\frac{11}{20}}=\GTenHLProbability{0.55}.
\]
The cooperative pays a \GTenHLValue{handling fee} below \GTenHLValue{7 USD} for a parcel with probability $\GTenHLProbability{55\%}$.

\textbf{(ii) At least \GTenHLValue{16 USD}.} Include both $\GTenHLValue{16}$ and $\GTenHLValue{23}$; the boundary mass at \GTenHLValue{16 USD} is included:
\[
P(\GTenHLValue{X}\geq\GTenHLValue{16})=\GTenHLProbability{p_X(16)}+\GTenHLProbability{p_X(23)}=\frac{\GTenHLProbability{2}+\GTenHLProbability{2}}{\GTenHLProbability{20}}=\GTenHLProbability{\frac15}=\GTenHLProbability{0.20}.
\]
The cooperative pays a \GTenHLValue{handling fee} of \GTenHLValue{16 USD} or more for a parcel with probability $\GTenHLProbability{20\%}$.

\par\smallskip
\noindent\begin{minipage}{\linewidth}\centering
\textbf{Proposed partial credit.}\par\smallskip
\begin{tabular}{@{}l r@{}}\toprule
\CreditDescription{\textbf{Assessed step or criterion}} & \textbf{Points}\\\midrule
\CreditDescription{(i): selecting exactly the included PMF masses with the correct boundary} & 0.125\\
\CreditDescription{(i): summing them correctly} & 0.0625\\
\CreditDescription{(i): the contextual interpretation} & 0.0625\\
\CreditDescription{(ii): selecting exactly the included PMF masses with the correct boundary} & 0.125\\
\CreditDescription{(ii): summing them correctly} & 0.0625\\
\CreditDescription{(ii): the contextual interpretation} & 0.0625\\
\midrule
\CreditDescription{\textbf{Total}} & \textbf{0.5}\\
\bottomrule\end{tabular}
\end{minipage}\par\smallskip

Each part totals 0.25 points.

\par\vspace{3mm}\secInicio[][sec:q3]

\textbf{Question total: $0.5+1+0.5=2$ points for C2.}

\clearpage
\sectionScored{Warranty service payouts}{C2-\sml{2}}{sec:q4}

\textbf{Basic Type A.} A warranty provider classifies one claim in exactly one of eight exhaustive, mutually exclusive categories: \GTenHLOutcome{no service}, \GTenHLOutcome{remote diagnosis}, \GTenHLOutcome{software reset}, \GTenHLOutcome{minor repair}, \GTenHLOutcome{major repair}, \GTenHLOutcome{replacement}, \GTenHLOutcome{store credit} or \GTenHLOutcome{cash refund}. These eight named categories are the distinct outcomes in the sample space. Their probabilities are \GTenHLProbability{0.20}, \GTenHLProbability{0.10}, \GTenHLProbability{0.15}, \GTenHLProbability{0.15}, \GTenHLProbability{0.10}, \GTenHLProbability{0.15}, \GTenHLProbability{0.05} and \GTenHLProbability{0.10}, respectively; do not assume equal likelihood. The corresponding payouts are \GTenHLValue{0}, \GTenHLValue{2}, \GTenHLValue{4}, \GTenHLValue{6}, \GTenHLValue{10}, \GTenHLValue{15}, \GTenHLValue{15} and \GTenHLValue{25 USD}, respectively. Count \GTenHLOutcome{store credit} at its face value as a payout. Each payout is fixed for its category, and no additional payments apply. Record the provider's payout for one claim.

\Needspace{8\baselineskip}
\subsection*{Task 1 --- question and worked solution (proposed allocation) \score{0.5}}
\label{sol:q4t1}

\noindent\textbf{\boldmath Define the \GTenHLValue{random variable} and \GTenHLValue{support}; derive and verify the grouped \GTenHLProbability{PMF}; produce its table and a labeled graph with actual monetary spacing.}\par\nopagebreak

Let $\GTenHLValue{X}$ be the provider payout for one claim, in \GTenHLValue{USD}. Map the supplied categories to payments:

\begin{center}
\begin{tabular}{ccc}\toprule
Category & Given probability & $\GTenHLValue{X}$ (\GTenHLValue{USD})\\\midrule
\GTenHLOutcome{No service} & $\GTenHLProbability{0.20}=\GTenHLProbability{4/20}$ & $\GTenHLValue{0}$\\
\GTenHLOutcome{Remote diagnosis} & $\GTenHLProbability{0.10}=\GTenHLProbability{2/20}$ & $\GTenHLValue{2}$\\
\GTenHLOutcome{Software reset} & $\GTenHLProbability{0.15}=\GTenHLProbability{3/20}$ & $\GTenHLValue{4}$\\
\GTenHLOutcome{Minor repair} & $\GTenHLProbability{0.15}=\GTenHLProbability{3/20}$ & $\GTenHLValue{6}$\\
\GTenHLOutcome{Major repair} & $\GTenHLProbability{0.10}=\GTenHLProbability{2/20}$ & $\GTenHLValue{10}$\\
\GTenHLOutcome{Replacement} & $\GTenHLProbability{0.15}=\GTenHLProbability{3/20}$ & $\GTenHLValue{15}$\\
\GTenHLOutcome{Store credit} & $\GTenHLProbability{0.05}=\GTenHLProbability{1/20}$ & $\GTenHLValue{15}$\\
\GTenHLOutcome{Cash refund} & $\GTenHLProbability{0.10}=\GTenHLProbability{2/20}$ & $\GTenHLValue{25}$\\
\bottomrule\end{tabular}
\end{center}

Replacement and store credit are distinct outcomes sharing a \GTenHLValue{15 USD} payout, so $\GTenHLProbability{p_X(15)}=\GTenHLProbability{0.15}+\GTenHLProbability{0.05}=\GTenHLProbability{0.20}=\GTenHLProbability{4/20}$. Store credit counts at face value. Every other category has its own payout and retains its supplied probability.

$\GTenHLValue{S_X}=\{\GTenHLValue{0},\GTenHLValue{2},\GTenHLValue{4},\GTenHLValue{6},\GTenHLValue{10},\GTenHLValue{15},\GTenHLValue{25}\}$ \GTenHLValue{USD}. There are eight elementary outcomes but seven distinct monetary values. Add probabilities of mutually exclusive categories sharing a value, because the events cannot occur together.

\begin{center}
\begin{minipage}{.26\linewidth}\centering
\begin{tabular}{cc}\toprule
$\GTenHLValue{x}$ (\GTenHLValue{USD}) & $\GTenHLProbability{p_X(x)}$\\\midrule
$\GTenHLValue{0}$ & $\GTenHLProbability{4/20}$\\
$\GTenHLValue{2}$ & $\GTenHLProbability{2/20}$\\
$\GTenHLValue{4}$ & $\GTenHLProbability{3/20}$\\
$\GTenHLValue{6}$ & $\GTenHLProbability{3/20}$\\
$\GTenHLValue{10}$ & $\GTenHLProbability{2/20}$\\
$\GTenHLValue{15}$ & $\GTenHLProbability{4/20}$\\
$\GTenHLValue{25}$ & $\GTenHLProbability{2/20}$\\
\bottomrule\end{tabular}
\end{minipage}\hfill
\begin{minipage}{.70\linewidth}\centering
\begin{tikzpicture}
\begin{axis}[keyaxis,width=.94\linewidth,height=5cm,xmin=-2,xmax=27,ymin=0,ymax=.35,xtick={0,2,4,6,10,15,25},xticklabel style={rotate=45,anchor=north east},ylabel={$\GTenHLProbability{p_X(x)}$},yticklabel={$\GTenHLProbability{\pgfmathprintnumber{\tick}}$},ytick={0,.05,.1,.15,.2,.25,.3}]
\draw[fill=keyMassFill,draw=keyMassColor] (axis cs:-0.45,0) rectangle (axis cs:0.45,4/20);
\draw[fill=keyMassFill,draw=keyMassColor] (axis cs:1.55,0) rectangle (axis cs:2.45,2/20);
\draw[fill=keyMassFill,draw=keyMassColor] (axis cs:3.55,0) rectangle (axis cs:4.45,3/20);
\draw[fill=keyMassFill,draw=keyMassColor] (axis cs:5.55,0) rectangle (axis cs:6.45,3/20);
\draw[fill=keyMassFill,draw=keyMassColor] (axis cs:9.55,0) rectangle (axis cs:10.45,2/20);
\draw[fill=keyMassFill,draw=keyMassColor] (axis cs:14.55,0) rectangle (axis cs:15.45,4/20);
\draw[fill=keyMassFill,draw=keyMassColor] (axis cs:24.55,0) rectangle (axis cs:25.45,2/20);
\end{axis}\end{tikzpicture}
\end{minipage}
\end{center}

\Needspace{5\baselineskip}\textbf{PMF checks.}
\begin{itemize}
\item Nonnegativity: every listed mass is strictly positive.
\item Normalization: $\GTenHLProbability{4}+\GTenHLProbability{2}+\GTenHLProbability{3}+\GTenHLProbability{3}+\GTenHLProbability{2}+\GTenHLProbability{4}+\GTenHLProbability{2}=\GTenHLProbability{20}$, so the masses sum to $\GTenHLProbability{20/20}=\GTenHLProbability{1}$.
\item Outside the \GTenHLValue{support}: set $\GTenHLProbability{p_X(x)}=\GTenHLProbability{0}$ there. The graph places bars at the actual monetary values; gaps carry no mass.
\end{itemize}

\textbf{Scoring note.} These explanatory positivity and normalization checks carry no separate points; omitting them does not reduce the task score.

\par\smallskip
\noindent\begin{minipage}{\linewidth}\centering
\textbf{Proposed partial credit.}\par\smallskip
\begin{tabular}{@{}l r@{}}\toprule
\CreditDescription{\textbf{Assessed step or criterion}} & \textbf{Points}\\\midrule
\CreditDescription{Variable, units and support} & 0.125\\
\CreditDescription{Category mapping, grouping and the PMF table} & 0.25\\
\CreditDescription{The labeled, correctly spaced graph} & 0.125\\
\midrule
\CreditDescription{\textbf{Total}} & \textbf{0.5}\\
\bottomrule\end{tabular}
\end{minipage}\par\smallskip

\par\vspace{3mm}\secInicio[][sec:q4]

\clearpage
\textbf{Question 4 --- continued}\quad\hyperref[sec:q4]{Original description}

\Needspace{8\baselineskip}
\subsection*{Task 2 --- question and worked solution (proposed allocation) \score{1}}
\label{sol:q4t2}

\noindent\textbf{\boldmath Construct the complete \HLCumulativeProbability{CDF} table, including both tails, and draw a labeled step graph. Use open and filled endpoints correctly and preserve actual monetary spacing. Verify the defining \HLCumulativeProbability{CDF} properties and that its jumps equal the \GTenHLProbability{PMF} masses.}\par\nopagebreak

For $\HLCumulativeProbability{F_X(x)}=P(\GTenHLValue{X}\leq \GTenHLValue{x})$, add masses in increasing payout order. The successive cumulative numerators are $\HLCumulativeProbability{4},\HLCumulativeProbability{6},\HLCumulativeProbability{9},\HLCumulativeProbability{12},\HLCumulativeProbability{14},\HLCumulativeProbability{18},\HLCumulativeProbability{20}$, all over $\HLCumulativeProbability{20}$.

\begin{center}
\begin{tabular}{cc}\toprule
Cutoff interval (\GTenHLValue{USD}) & $\HLCumulativeProbability{F_X(x)}$\\\midrule
$\GTenHLValue{x}<\GTenHLValue{0}$ & $\HLCumulativeProbability{0}$\\
$\GTenHLValue{0}\leq \GTenHLValue{x}<\GTenHLValue{2}$ & $\HLCumulativeProbability{4/20}$\\
$\GTenHLValue{2}\leq \GTenHLValue{x}<\GTenHLValue{4}$ & $\HLCumulativeProbability{6/20}$\\
$\GTenHLValue{4}\leq \GTenHLValue{x}<\GTenHLValue{6}$ & $\HLCumulativeProbability{9/20}$\\
$\GTenHLValue{6}\leq \GTenHLValue{x}<\GTenHLValue{10}$ & $\HLCumulativeProbability{12/20}$\\
$\GTenHLValue{10}\leq \GTenHLValue{x}<\GTenHLValue{15}$ & $\HLCumulativeProbability{14/20}$\\
$\GTenHLValue{15}\leq \GTenHLValue{x}<\GTenHLValue{25}$ & $\HLCumulativeProbability{18/20}$\\
$\GTenHLValue{x}\geq \GTenHLValue{25}$ & $\HLCumulativeProbability{1}$\\
\bottomrule\end{tabular}
\end{center}

\begin{center}\begin{tikzpicture}
\begin{axis}[keyaxis,width=.86\linewidth,height=7cm,xmin=-2,xmax=27,ymin=-.07,ymax=1.1,xtick={0,2,4,6,10,15,25},xticklabel style={rotate=45,anchor=north east},ytick={0,.25,.5,.75,1},ylabel={$\HLCumulativeProbability{F_X(x)}$},yticklabel={$\HLCumulativeProbability{\pgfmathprintnumber{\tick}}$}]
\addplot[keyCDFColor,thick,<-] coordinates {(-2,0)(0,0)};
\addplot[keyCDFColor,thick] coordinates {(0,4/20)(2,4/20)};
\addplot[keyCDFColor,thick] coordinates {(2,6/20)(4,6/20)};
\addplot[keyCDFColor,thick] coordinates {(4,9/20)(6,9/20)};
\addplot[keyCDFColor,thick] coordinates {(6,12/20)(10,12/20)};
\addplot[keyCDFColor,thick] coordinates {(10,14/20)(15,14/20)};
\addplot[keyCDFColor,thick] coordinates {(15,18/20)(25,18/20)};
\addplot[keyCDFColor,thick,->] coordinates {(25,1)(27,1)};
\addplot[keyCDFColor,only marks,mark=*,mark size=1.6pt,mark options={fill=white}] coordinates {(0,0/20)(2,4/20)(4,6/20)(6,9/20)(10,12/20)(15,14/20)(25,18/20)};
\addplot[keyCDFColor,only marks,mark=*,mark size=1.6pt] coordinates {(0,4/20)(2,6/20)(4,9/20)(6,12/20)(10,14/20)(15,18/20)(25,20/20)};
\end{axis}\end{tikzpicture}\end{center}

\textbf{CDF checks.}
\begin{itemize}
\item Domain: the table covers every real cutoff.
\item Range: all tabulated heights lie in $[\HLCumulativeProbability{0},\HLCumulativeProbability{1}]$.
\item Monotonicity: the tabulated heights never decrease.
\item Right-continuity: each horizontal segment includes its left endpoint and excludes its right endpoint; the filled point gives the value at a jump.
\item Lower-tail limit: the left arrow extends to $-\infty$, where the limit is 0.
\item Upper-tail limit: the right arrow extends to $+\infty$, where the limit is 1.
\item Jump--mass agreement: in \GTenHLValue{support} order, subtracting successive cumulative numerators (starting at zero) gives $\GTenHLProbability{4},\GTenHLProbability{2},\GTenHLProbability{3},\GTenHLProbability{3},\GTenHLProbability{2},\GTenHLProbability{4},\GTenHLProbability{2}$. Hence each jump is its \GTenHLProbability{PMF} mass over $\GTenHLProbability{20}$: $\HLCumulativeProbability{F_X(a)}-\HLCumulativeProbability{F_X(a^-)}=\GTenHLProbability{p_X(a)}$. At unsupported cutoffs both the jump and the mass are zero.
\end{itemize}

\par\smallskip
\noindent\begin{minipage}{\linewidth}\centering
\textbf{Proposed partial credit.}\par\smallskip
\begin{tabular}{@{}l r@{}}\toprule
\CreditDescription{\textbf{Assessed step or criterion}} & \textbf{Points}\\\midrule
\CreditDescription{Cumulative sums} & 0.25\\
\CreditDescription{The complete interval table and tails} & 0.25\\
\CreditDescription{The labeled, correctly spaced step graph and endpoints} & 0.25\\
\CreditDescription{The CDF and jump checks} & 0.25\\
\midrule
\CreditDescription{\textbf{Total}} & \textbf{1}\\
\bottomrule\end{tabular}
\end{minipage}\par\smallskip

\par\vspace{3mm}\secInicio[][sec:q4]

\clearpage
\textbf{Question 4 --- continued}\quad\hyperref[sec:q4]{Original description}

\Needspace{8\baselineskip}
\subsection*{Task 3 (c) --- question and worked solution (proposed allocation) \score{0.5}}
\label{sol:q4t3}

\noindent\textbf{\boldmath Calculate and interpret (i) the probability that the provider's payout is less than \GTenHLValue{6 USD} and (ii) the probability that it is at least \GTenHLValue{15 USD}.}\par\nopagebreak

\textbf{Allocation:} (i) 0.25 points; (ii) 0.25 points.

\textbf{(i) Less than \GTenHLValue{6 USD}.} Include only the \GTenHLValue{support} values $\GTenHLValue{0},\GTenHLValue{2},\GTenHLValue{4}$; the mass at \GTenHLValue{6 USD} is excluded because the inequality is strict:
\[
P(\GTenHLValue{X}<\GTenHLValue{6})=\GTenHLProbability{p_X(0)}+\GTenHLProbability{p_X(2)}+\GTenHLProbability{p_X(4)}
=\frac{\GTenHLProbability{4}+\GTenHLProbability{2}+\GTenHLProbability{3}}{\GTenHLProbability{20}}=\GTenHLProbability{\frac9{20}}=\GTenHLProbability{0.45}.
\]
The provider pays less than \GTenHLValue{6 USD} for one claim with probability $\GTenHLProbability{45\%}$.

\textbf{(ii) At least \GTenHLValue{15 USD}.} Include both $\GTenHLValue{15}$ and $\GTenHLValue{25}$; the boundary mass at \GTenHLValue{15 USD} includes replacement and store credit:
\[
P(\GTenHLValue{X}\geq\GTenHLValue{15})=\GTenHLProbability{p_X(15)}+\GTenHLProbability{p_X(25)}=\frac{\GTenHLProbability{4}+\GTenHLProbability{2}}{\GTenHLProbability{20}}=\GTenHLProbability{\frac3{10}}=\GTenHLProbability{0.30}.
\]
The provider pays \GTenHLValue{15 USD} or more for one claim with probability $\GTenHLProbability{30\%}$, counting store credit at face value.

\par\smallskip
\noindent\begin{minipage}{\linewidth}\centering
\textbf{Proposed partial credit.}\par\smallskip
\begin{tabular}{@{}l r@{}}\toprule
\CreditDescription{\textbf{Assessed step or criterion}} & \textbf{Points}\\\midrule
\CreditDescription{(i): selecting exactly the included PMF masses with the correct boundary} & 0.125\\
\CreditDescription{(i): summing them correctly} & 0.0625\\
\CreditDescription{(i): the contextual interpretation} & 0.0625\\
\CreditDescription{(ii): selecting exactly the included PMF masses with the correct boundary} & 0.125\\
\CreditDescription{(ii): summing them correctly} & 0.0625\\
\CreditDescription{(ii): the contextual interpretation} & 0.0625\\
\midrule
\CreditDescription{\textbf{Total}} & \textbf{0.5}\\
\bottomrule\end{tabular}
\end{minipage}\par\smallskip

Each part totals 0.25 points.

\par\vspace{3mm}\secInicio[][sec:q4]

\textbf{Question total: $0.5+1+0.5=2$ points for C2.}

\clearpage
\sectionScored{Marketplace account credits}{C2-\sml{4}}{sec:q5}

\textbf{Advanced Type A.} For one purchase, a marketplace awards exactly one \GTenHLValue{account credit}: \GTenHLValue{0}, \GTenHLValue{2}, \GTenHLValue{5}, \GTenHLValue{9}, \GTenHLValue{14} or \GTenHLValue{20 USD}. These are all possible credits and their categories are mutually exclusive. A \GTenHLValue{2 USD} credit is twice as likely as no credit. A \GTenHLValue{5 USD} credit is as likely as no credit and a \GTenHLValue{2 USD} credit combined. A \GTenHLValue{9 USD} credit is twice as likely as a \GTenHLValue{5 USD} credit. A \GTenHLValue{14 USD} credit is half as likely as a \GTenHLValue{9 USD} credit. A \GTenHLValue{20 USD} credit is as likely as the \GTenHLValue{2 USD} and \GTenHLValue{14 USD} credits combined. There are no other payments or adjustments. Determine the probabilities from these likelihood relationships and record the \GTenHLValue{account credit}.

\Needspace{8\baselineskip}
\subsection*{Task 1 --- question and worked solution (proposed allocation) \score{1}}
\label{sol:q5t1}

\noindent\textbf{\boldmath Define the \GTenHLValue{random variable} and \GTenHLValue{support}. Derive the six probabilities using normalization, verify the \GTenHLProbability{PMF}, and produce its table and labeled graph.}\par\nopagebreak

Let $\GTenHLValue{X}$ be the \GTenHLValue{account credit} for one purchase, in \GTenHLValue{USD}, with $\GTenHLValue{S_X}=\{\GTenHLValue{0},\GTenHLValue{2},\GTenHLValue{5},\GTenHLValue{9},\GTenHLValue{14},\GTenHLValue{20}\}$. Set $\GTenHLParameter{a}=P(\GTenHLValue{X}=\GTenHLValue{0})$. Since the categories are mutually exclusive, a combined likelihood is a sum. Translate the relationships in their stated order:
\begin{align*}
\GTenHLProbability{p_X(0)}&=\GTenHLParameter{a},&\GTenHLProbability{p_X(2)}&=2\GTenHLParameter{a},\\
\GTenHLProbability{p_X(5)}&=\GTenHLParameter{a}+2\GTenHLParameter{a}=3\GTenHLParameter{a},&\GTenHLProbability{p_X(9)}&=2(3\GTenHLParameter{a})=6\GTenHLParameter{a},\\
\GTenHLProbability{p_X(14)}&=\tfrac12(6\GTenHLParameter{a})=3\GTenHLParameter{a},&\GTenHLProbability{p_X(20)}&=2\GTenHLParameter{a}+3\GTenHLParameter{a}=5\GTenHLParameter{a}.
\end{align*}
Normalization gives
\[
\GTenHLParameter{a}+2\GTenHLParameter{a}+3\GTenHLParameter{a}+6\GTenHLParameter{a}+3\GTenHLParameter{a}+5\GTenHLParameter{a}=\GTenHLProbability{1},\qquad 20\GTenHLParameter{a}=\GTenHLProbability{1},\qquad \GTenHLParameter{a}=\GTenHLParameter{\frac1{20}}.
\]
Every probability is now uniquely determined: all depend on the single baseline, whose positive coefficient in the normalization equation is 20. No probability is zero. Each credit is already a distinct monetary value, so no additional grouping is needed.

\begin{center}
\begin{minipage}{.26\linewidth}\centering
\begin{tabular}{cc}\toprule
$\GTenHLValue{x}$ (\GTenHLValue{USD}) & $\GTenHLProbability{p_X(x)}$\\\midrule
$\GTenHLValue{0}$ & $\GTenHLProbability{1/20}$\\
$\GTenHLValue{2}$ & $\GTenHLProbability{2/20}=\GTenHLProbability{1/10}$\\
$\GTenHLValue{5}$ & $\GTenHLProbability{3/20}$\\
$\GTenHLValue{9}$ & $\GTenHLProbability{6/20}=\GTenHLProbability{3/10}$\\
$\GTenHLValue{14}$ & $\GTenHLProbability{3/20}$\\
$\GTenHLValue{20}$ & $\GTenHLProbability{5/20}=\GTenHLProbability{1/4}$\\
\bottomrule\end{tabular}
\end{minipage}\hfill
\begin{minipage}{.70\linewidth}\centering
\begin{tikzpicture}
\begin{axis}[keyaxis,width=.94\linewidth,height=5cm,xmin=-2,xmax=22,ymin=0,ymax=.65,xtick={0,2,5,9,14,20},xticklabel style={rotate=45,anchor=north east},ylabel={$\GTenHLProbability{p_X(x)}$},yticklabel={$\GTenHLProbability{\pgfmathprintnumber{\tick}}$},ytick={0,.1,.2,.3,.4,.5,.6}]

\draw[fill=keyMassFill,draw=keyMassColor] (axis cs:-0.45,0) rectangle (axis cs:0.45,1/20);

\draw[fill=keyMassFill,draw=keyMassColor] (axis cs:1.55,0) rectangle (axis cs:2.45,2/20);

\draw[fill=keyMassFill,draw=keyMassColor] (axis cs:4.55,0) rectangle (axis cs:5.45,3/20);

\draw[fill=keyMassFill,draw=keyMassColor] (axis cs:8.55,0) rectangle (axis cs:9.45,6/20);

\draw[fill=keyMassFill,draw=keyMassColor] (axis cs:13.55,0) rectangle (axis cs:14.45,3/20);

\draw[fill=keyMassFill,draw=keyMassColor] (axis cs:19.55,0) rectangle (axis cs:20.45,5/20);

\end{axis}\end{tikzpicture}
\end{minipage}
\end{center}

\Needspace{5\baselineskip}\textbf{PMF checks.}
\begin{itemize}
\item Nonnegativity: every listed mass is strictly positive.
\item Normalization: $\GTenHLProbability{1}+\GTenHLProbability{2}+\GTenHLProbability{3}+\GTenHLProbability{6}+\GTenHLProbability{3}+\GTenHLProbability{5}=\GTenHLProbability{20}$, so the masses sum to $\GTenHLProbability{20/20}=\GTenHLProbability{1}$.
\item Outside the \GTenHLValue{support}: set $\GTenHLProbability{p_X(x)}=\GTenHLProbability{0}$ there. The graph places bars at the actual monetary values; gaps carry no mass.
\end{itemize}

\textbf{Relationship check.}
\begin{itemize}
\item Likelihood relationships: the numerators satisfy $\GTenHLProbability{2}=2(\GTenHLProbability{1})$, $\GTenHLProbability{3}=\GTenHLProbability{1}+\GTenHLProbability{2}$, $\GTenHLProbability{6}=2(\GTenHLProbability{3})$, $\GTenHLProbability{3}=\GTenHLProbability{6/2}$, and $\GTenHLProbability{5}=\GTenHLProbability{2}+\GTenHLProbability{3}$, exactly reproducing all five verbal relationships.
\end{itemize}

\textbf{Scoring note.} These explanatory positivity and normalization checks carry no separate points; omitting them does not reduce the task score.

\par\smallskip
\noindent\begin{minipage}{\linewidth}\centering
\textbf{Proposed partial credit.}\par\smallskip
\begin{tabular}{@{}l r@{}}\toprule
\CreditDescription{\textbf{Assessed step or criterion}} & \textbf{Points}\\\midrule
\CreditDescription{Variable, units and support} & 0.125\\
\CreditDescription{Translating the relationships and solving normalization} & 0.375\\
\CreditDescription{All six PMF entries} & 0.25\\
\CreditDescription{The labeled, correctly spaced graph} & 0.125\\
\CreditDescription{Relationship checks} & 0.125\\
\midrule
\CreditDescription{\textbf{Total}} & \textbf{1}\\
\bottomrule\end{tabular}
\end{minipage}\par\smallskip

\par\vspace{3mm}\secInicio[][sec:q5]

\clearpage
\textbf{Question 5 --- continued}\quad\hyperref[sec:q5]{Original description}

\Needspace{8\baselineskip}
\subsection*{Task 2 --- question and worked solution (proposed allocation) \score{2}}
\label{sol:q5t2}

\noindent\textbf{\boldmath Construct the complete \HLCumulativeProbability{CDF} table, including both tails, and draw a labeled step graph. Use open and filled endpoints correctly and preserve actual monetary spacing. Verify the defining \HLCumulativeProbability{CDF} properties and that its jumps equal the \GTenHLProbability{PMF} masses.}\par\nopagebreak

For $\HLCumulativeProbability{F_X(x)}=P(\GTenHLValue{X}\leq \GTenHLValue{x})$, add masses in increasing payout order. The successive cumulative numerators are $\HLCumulativeProbability{1},\HLCumulativeProbability{3},\HLCumulativeProbability{6},\HLCumulativeProbability{12},\HLCumulativeProbability{15},\HLCumulativeProbability{20}$, all over $\HLCumulativeProbability{20}$.

\begin{center}
\begin{tabular}{cc}\toprule
Cutoff interval (\GTenHLValue{USD}) & $\HLCumulativeProbability{F_X(x)}$\\\midrule
$\GTenHLValue{x}<\GTenHLValue{0}$ & $\HLCumulativeProbability{0}$\\
$\GTenHLValue{0}\leq \GTenHLValue{x}<\GTenHLValue{2}$ & $\HLCumulativeProbability{1/20}$\\
$\GTenHLValue{2}\leq \GTenHLValue{x}<\GTenHLValue{5}$ & $\HLCumulativeProbability{3/20}$\\
$\GTenHLValue{5}\leq \GTenHLValue{x}<\GTenHLValue{9}$ & $\HLCumulativeProbability{3/10}$\\
$\GTenHLValue{9}\leq \GTenHLValue{x}<\GTenHLValue{14}$ & $\HLCumulativeProbability{3/5}$\\
$\GTenHLValue{14}\leq \GTenHLValue{x}<\GTenHLValue{20}$ & $\HLCumulativeProbability{3/4}$\\
$\GTenHLValue{x}\geq \GTenHLValue{20}$ & $\HLCumulativeProbability{1}$\\
\bottomrule\end{tabular}
\end{center}

\begin{center}\begin{tikzpicture}
\begin{axis}[keyaxis,width=.86\linewidth,height=7cm,xmin=-2,xmax=22,ymin=-.07,ymax=1.1,xtick={0,2,5,9,14,20},xticklabel style={rotate=45,anchor=north east},ytick={0,.25,.5,.75,1},ylabel={$\HLCumulativeProbability{F_X(x)}$},yticklabel={$\HLCumulativeProbability{\pgfmathprintnumber{\tick}}$}]

\addplot[keyCDFColor,thick,<-] coordinates {(-2,0)(0,0)};

\addplot[keyCDFColor,thick] coordinates {(0,1/20)(2,1/20)};

\addplot[keyCDFColor,thick] coordinates {(2,3/20)(5,3/20)};

\addplot[keyCDFColor,thick] coordinates {(5,6/20)(9,6/20)};

\addplot[keyCDFColor,thick] coordinates {(9,12/20)(14,12/20)};

\addplot[keyCDFColor,thick] coordinates {(14,15/20)(20,15/20)};

\addplot[keyCDFColor,thick,->] coordinates {(20,1)(22,1)};

\addplot[keyCDFColor,only marks,mark=*,mark size=1.6pt,mark options={fill=white}] coordinates {(0,0/20)(2,1/20)(5,3/20)(9,6/20)(14,12/20)(20,15/20)};

\addplot[keyCDFColor,only marks,mark=*,mark size=1.6pt] coordinates {(0,1/20)(2,3/20)(5,6/20)(9,12/20)(14,15/20)(20,20/20)};

\end{axis}\end{tikzpicture}\end{center}

\textbf{CDF checks.}
\begin{itemize}
\item Domain: the table covers every real cutoff.
\item Range: all tabulated heights lie in $[\HLCumulativeProbability{0},\HLCumulativeProbability{1}]$.
\item Monotonicity: the tabulated heights never decrease.
\item Right-continuity: each horizontal segment includes its left endpoint and excludes its right endpoint; the filled point gives the value at a jump.
\item Lower-tail limit: the left arrow extends to $-\infty$, where the limit is 0.
\item Upper-tail limit: the right arrow extends to $+\infty$, where the limit is 1.
\item Jump--mass agreement: in \GTenHLValue{support} order, subtracting successive cumulative numerators (starting at zero) gives $\GTenHLProbability{1},\GTenHLProbability{2},\GTenHLProbability{3},\GTenHLProbability{6},\GTenHLProbability{3},\GTenHLProbability{5}$. Hence each jump is its \GTenHLProbability{PMF} mass over $\GTenHLProbability{20}$: $\HLCumulativeProbability{F_X(a)}-\HLCumulativeProbability{F_X(a^-)}=\GTenHLProbability{p_X(a)}$. At unsupported cutoffs both the jump and the mass are zero.
\end{itemize}

\par\smallskip
\noindent\begin{minipage}{\linewidth}\centering
\textbf{Proposed partial credit.}\par\smallskip
\begin{tabular}{@{}l r@{}}\toprule
\CreditDescription{\textbf{Assessed step or criterion}} & \textbf{Points}\\\midrule
\CreditDescription{Cumulative sums} & 0.5\\
\CreditDescription{The complete interval table and tails} & 0.5\\
\CreditDescription{The labeled, correctly spaced step graph and endpoints} & 0.5\\
\CreditDescription{The CDF and jump checks} & 0.5\\
\midrule
\CreditDescription{\textbf{Total}} & \textbf{2}\\
\bottomrule\end{tabular}
\end{minipage}\par\smallskip

\par\vspace{3mm}\secInicio[][sec:q5]

\clearpage
\textbf{Question 5 --- continued}\quad\hyperref[sec:q5]{Original description}

\Needspace{8\baselineskip}
\subsection*{Task 3 (c) --- question and worked solution (proposed allocation) \score{1}}
\label{sol:q5t3}

\noindent\textbf{\boldmath Calculate and interpret (i) the probability that the \GTenHLValue{account credit} is less than \GTenHLValue{9 USD} and (ii) the probability that it is at least \GTenHLValue{14 USD}.}\par\nopagebreak

\textbf{Allocation:} (i) 0.5 points; (ii) 0.5 points.

\textbf{(i) Less than \GTenHLValue{9 USD}.} Include only the \GTenHLValue{support} values $\GTenHLValue{0},\GTenHLValue{2},\GTenHLValue{5}$; the mass at \GTenHLValue{9 USD} is excluded because the inequality is strict:
\[
P(\GTenHLValue{X}<\GTenHLValue{9})=\GTenHLProbability{p_X(0)}+\GTenHLProbability{p_X(2)}+\GTenHLProbability{p_X(5)}
=\frac{\GTenHLProbability{1}+\GTenHLProbability{2}+\GTenHLProbability{3}}{\GTenHLProbability{20}}=\GTenHLProbability{\frac3{10}}=\GTenHLProbability{0.30}.
\]
For one purchase, the \GTenHLValue{account credit} is below \GTenHLValue{9 USD} with probability $\GTenHLProbability{30\%}$.

\textbf{(ii) At least \GTenHLValue{14 USD}.} Include both $\GTenHLValue{14}$ and $\GTenHLValue{20}$; the boundary mass at \GTenHLValue{14 USD} is included:
\[
P(\GTenHLValue{X}\geq\GTenHLValue{14})=\GTenHLProbability{p_X(14)}+\GTenHLProbability{p_X(20)}=\frac{\GTenHLProbability{3}+\GTenHLProbability{5}}{\GTenHLProbability{20}}=\GTenHLProbability{\frac25}=\GTenHLProbability{0.40}.
\]
For one purchase, the \GTenHLValue{account credit} is \GTenHLValue{14 USD} or more with probability $\GTenHLProbability{40\%}$.

\par\smallskip
\noindent\begin{minipage}{\linewidth}\centering
\textbf{Proposed partial credit.}\par\smallskip
\begin{tabular}{@{}l r@{}}\toprule
\CreditDescription{\textbf{Assessed step or criterion}} & \textbf{Points}\\\midrule
\CreditDescription{(i): selecting exactly the included PMF masses with the correct boundary} & 0.25\\
\CreditDescription{(i): summing them correctly} & 0.125\\
\CreditDescription{(i): the contextual interpretation} & 0.125\\
\CreditDescription{(ii): selecting exactly the included PMF masses with the correct boundary} & 0.25\\
\CreditDescription{(ii): summing them correctly} & 0.125\\
\CreditDescription{(ii): the contextual interpretation} & 0.125\\
\midrule
\CreditDescription{\textbf{Total}} & \textbf{1}\\
\bottomrule\end{tabular}
\end{minipage}\par\smallskip

Each part totals 0.5 points.

\par\vspace{3mm}\secInicio[][sec:q5]

\textbf{Question total: $1+2+1=4$ points for C2.}

\clearpage
\sectionScored{Community market prize tickets}{C2-\sml{6}}{sec:q6}

\textbf{Advanced Type B.} A community market has nine physically distinct prize tickets: four worth \GTenHLValue{0 USD} each, three worth \GTenHLValue{5 USD} each and two worth \GTenHLValue{12 USD} each. Choose a subset of three tickets uniformly without \GTenHLOutcome{replacement}, so every subset of three physically distinct tickets has the same chance. Ignore selection order. First add the face values. Then add one bonus of \GTenHLValue{7 USD} if the selected subset contains at least one \GTenHLValue{5 USD} ticket and at least one \GTenHLValue{12 USD} ticket; otherwise add no bonus. Finally, if this adjusted amount exceeds \GTenHLValue{25 USD}, pay only \GTenHLValue{25 USD}; otherwise pay the adjusted amount. No other payments or deductions apply. Record the \GTenHLValue{final payout}.

\Needspace{8\baselineskip}
\subsection*{Task 1 --- question and worked solution (proposed allocation) \score{1.5}}
\label{sol:q6t1}

\noindent\textbf{\boldmath Define the \GTenHLValue{random variable} and \GTenHLValue{support}. Count the distinct-ticket selections, apply the bonus and then the cap, aggregate equal final payouts, and derive and verify the \GTenHLProbability{PMF}. Produce its table and labeled graph.}\par\nopagebreak

Let $\GTenHLValue{X}$ be the \GTenHLValue{final payout} for one three-ticket subset, in \GTenHLValue{USD}. Label the physical tickets $\GTenHLOutcome{Z_1},\ldots,\GTenHLOutcome{Z_4}$, $\GTenHLOutcome{L_1},\ldots,\GTenHLOutcome{L_3}$, and $\GTenHLOutcome{H_1},\GTenHLOutcome{H_2}$, with values 0, 5 and \GTenHLValue{12 USD}. There are
\[
\binom{9}{3}=\frac{9\cdot8\cdot7}{3\cdot2\cdot1}=84
\]
equally likely unordered subsets, each with probability $\GTenHLProbability{1/84}$. A composition records $(a,b,c)$, the counts of zero-, five- and twelve-dollar tickets. Feasible counts satisfy $a+b+c=3$, $0\leq a\leq4$, $0\leq b\leq3$, and $0\leq c\leq2$. The number of subsets with that composition is
\[
\GTenHLSecond{N}(a,b,c)=\binom{4}{a}\binom{3}{b}\binom{2}{c},\qquad P(a,b,c)=\frac{\GTenHLSecond{N}(a,b,c)}{\GTenHLProbability{84}}.
\]
Define $I=1$ if $b\geq1$ and $c\geq1$, and $I=0$ otherwise. The face sum is $\GTenHLValue{5}b+\GTenHLValue{12}c$; the adjusted amount is $\GTenHLSecond{Y}=\GTenHLValue{5}b+\GTenHLValue{12}c+\GTenHLValue{7}I$; finally $\GTenHLValue{X}=\min\{\GTenHLSecond{Y},\GTenHLValue{25}\}$. The bonus is added once, before the cap.

\begingroup\renewcommand{\arraystretch}{1.35}

\begin{center}
\begin{tabular}{ccccccc}\toprule
$(a,b,c)$ & Subset count & Probability & Face sum & Bonus & $\GTenHLSecond{Y}$ & $\GTenHLValue{X}$\\\midrule
$(\GTenHLOutcome{3},\GTenHLOutcome{0},\GTenHLOutcome{0})$ & $\binom4{3}\binom3{0}\binom2{0}=4$ & $\GTenHLProbability{4/84}$ & $\GTenHLValue{0}$ & $\GTenHLValue{0}$ & $\GTenHLSecond{0}$ & $\GTenHLValue{0}$\\
$(\GTenHLOutcome{2},\GTenHLOutcome{1},\GTenHLOutcome{0})$ & $\binom4{2}\binom3{1}\binom2{0}=18$ & $\GTenHLProbability{18/84}$ & $\GTenHLValue{5}$ & $\GTenHLValue{0}$ & $\GTenHLSecond{5}$ & $\GTenHLValue{5}$\\
$(\GTenHLOutcome{1},\GTenHLOutcome{2},\GTenHLOutcome{0})$ & $\binom4{1}\binom3{2}\binom2{0}=12$ & $\GTenHLProbability{12/84}$ & $\GTenHLValue{10}$ & $\GTenHLValue{0}$ & $\GTenHLSecond{10}$ & $\GTenHLValue{10}$\\
$(\GTenHLOutcome{0},\GTenHLOutcome{3},\GTenHLOutcome{0})$ & $\binom4{0}\binom3{3}\binom2{0}=1$ & $\GTenHLProbability{1/84}$ & $\GTenHLValue{15}$ & $\GTenHLValue{0}$ & $\GTenHLSecond{15}$ & $\GTenHLValue{15}$\\
$(\GTenHLOutcome{2},\GTenHLOutcome{0},\GTenHLOutcome{1})$ & $\binom4{2}\binom3{0}\binom2{1}=12$ & $\GTenHLProbability{12/84}$ & $\GTenHLValue{12}$ & $\GTenHLValue{0}$ & $\GTenHLSecond{12}$ & $\GTenHLValue{12}$\\
$(\GTenHLOutcome{1},\GTenHLOutcome{1},\GTenHLOutcome{1})$ & $\binom4{1}\binom3{1}\binom2{1}=24$ & $\GTenHLProbability{24/84}$ & $\GTenHLValue{17}$ & $\GTenHLValue{7}$ & $\GTenHLSecond{24}$ & $\GTenHLValue{24}$\\
$(\GTenHLOutcome{0},\GTenHLOutcome{2},\GTenHLOutcome{1})$ & $\binom4{0}\binom3{2}\binom2{1}=6$ & $\GTenHLProbability{6/84}$ & $\GTenHLValue{22}$ & $\GTenHLValue{7}$ & $\GTenHLSecond{29}$ & $\GTenHLValue{25}$\\
$(\GTenHLOutcome{1},\GTenHLOutcome{0},\GTenHLOutcome{2})$ & $\binom4{1}\binom3{0}\binom2{2}=4$ & $\GTenHLProbability{4/84}$ & $\GTenHLValue{24}$ & $\GTenHLValue{0}$ & $\GTenHLSecond{24}$ & $\GTenHLValue{24}$\\
$(\GTenHLOutcome{0},\GTenHLOutcome{1},\GTenHLOutcome{2})$ & $\binom4{0}\binom3{1}\binom2{2}=3$ & $\GTenHLProbability{3/84}$ & $\GTenHLValue{29}$ & $\GTenHLValue{7}$ & $\GTenHLSecond{36}$ & $\GTenHLValue{25}$\\
\bottomrule\end{tabular}
\end{center}

\endgroup

All monetary columns are in USD. These nine compositions are exhaustive: list $c=0,1,2$ and every feasible $b$, then $a=3-b-c$. Their counts sum to $4+18+12+1+12+24+6+4+3=84$. Compositions are not equally likely: for example, $(1,1,1)$ represents 24 subsets, whereas $(0,3,0)$ represents only one.

\textbf{Aggregation.} The \GTenHLValue{24 USD} payout comes from $(\GTenHLOutcome{1},\GTenHLOutcome{1},\GTenHLOutcome{1})$ (24 subsets after the bonus) and $(\GTenHLOutcome{1},\GTenHLOutcome{0},\GTenHLOutcome{2})$ (4 subsets without a bonus), so $\GTenHLProbability{p_X(24)}=(\GTenHLProbability{24}+\GTenHLProbability{4})/\GTenHLProbability{84}=\GTenHLProbability{28/84}$. The cap merges the adjusted amounts 29 and \GTenHLValue{36 USD} from $(\GTenHLOutcome{0},\GTenHLOutcome{2},\GTenHLOutcome{1})$ and $(\GTenHLOutcome{0},\GTenHLOutcome{1},\GTenHLOutcome{2})$, so $\GTenHLProbability{p_X(25)}=(\GTenHLProbability{6}+\GTenHLProbability{3})/\GTenHLProbability{84}=\GTenHLProbability{9/84}$. Each other payout arises from one composition.

Thus $\GTenHLValue{S_X}=\{\GTenHLValue{0},\GTenHLValue{5},\GTenHLValue{10},\GTenHLValue{12},\GTenHLValue{15},\GTenHLValue{24},\GTenHLValue{25}\}$ \GTenHLValue{USD}. The cap binds only when the adjusted amount is greater than \GTenHLValue{25 USD}; the \GTenHLValue{24 USD} payout remains unchanged.

\par\vspace{3mm}\secInicio[][sec:q6]

\clearpage\textbf{Question 6 --- Task 1 continued: grouped \GTenHLProbability{PMF} and graph}

\begin{center}
\begin{tabular}{cc}\toprule
$\GTenHLValue{x}$ (\GTenHLValue{USD}) & $\GTenHLProbability{p_X(x)}$\\\midrule
$\GTenHLValue{0}$ & $\GTenHLProbability{4/84}=\GTenHLProbability{1/21}$\\
$\GTenHLValue{5}$ & $\GTenHLProbability{18/84}=\GTenHLProbability{3/14}$\\
$\GTenHLValue{10}$ & $\GTenHLProbability{12/84}=\GTenHLProbability{1/7}$\\
$\GTenHLValue{12}$ & $\GTenHLProbability{12/84}=\GTenHLProbability{1/7}$\\
$\GTenHLValue{15}$ & $\GTenHLProbability{1/84}$\\
$\GTenHLValue{24}$ & $\GTenHLProbability{28/84}=\GTenHLProbability{1/3}$\\
$\GTenHLValue{25}$ & $\GTenHLProbability{9/84}=\GTenHLProbability{3/28}$\\
\bottomrule\end{tabular}
\end{center}

\begin{center}\begin{tikzpicture}
\begin{axis}[keyaxis,width=.86\linewidth,height=5cm,xmin=-2,xmax=27,ymin=0,ymax=.65,xtick={0,5,10,12,15,24,25},xticklabel style={rotate=45,anchor=north east},ylabel={$\GTenHLProbability{p_X(x)}$},yticklabel={$\GTenHLProbability{\pgfmathprintnumber{\tick}}$},ytick={0,.1,.2,.3,.4,.5,.6}]

\draw[fill=keyMassFill,draw=keyMassColor] (axis cs:-0.225,0) rectangle (axis cs:0.225,4/84);

\draw[fill=keyMassFill,draw=keyMassColor] (axis cs:4.775,0) rectangle (axis cs:5.225,18/84);

\draw[fill=keyMassFill,draw=keyMassColor] (axis cs:9.775,0) rectangle (axis cs:10.225,12/84);

\draw[fill=keyMassFill,draw=keyMassColor] (axis cs:11.775,0) rectangle (axis cs:12.225,12/84);

\draw[fill=keyMassFill,draw=keyMassColor] (axis cs:14.775,0) rectangle (axis cs:15.225,1/84);

\draw[fill=keyMassFill,draw=keyMassColor] (axis cs:23.775,0) rectangle (axis cs:24.225,28/84);

\draw[fill=keyMassFill,draw=keyMassColor] (axis cs:24.775,0) rectangle (axis cs:25.225,9/84);

\end{axis}\end{tikzpicture}\end{center}

\Needspace{5\baselineskip}\textbf{PMF checks.}
\begin{itemize}
\item Nonnegativity: every listed mass is strictly positive.
\item Normalization: $\GTenHLProbability{4}+\GTenHLProbability{18}+\GTenHLProbability{12}+\GTenHLProbability{12}+\GTenHLProbability{1}+\GTenHLProbability{28}+\GTenHLProbability{9}=\GTenHLProbability{84}$, so the masses sum to $\GTenHLProbability{84/84}=\GTenHLProbability{1}$.
\item Outside the \GTenHLValue{support}: set $\GTenHLProbability{p_X(x)}=\GTenHLProbability{0}$ there. The graph places bars at the actual monetary values; gaps carry no mass.
\end{itemize}

\textbf{Scoring note.} These explanatory positivity and normalization checks carry no separate points; omitting them does not reduce the task score.

\par\smallskip
\noindent\begin{minipage}{\linewidth}\centering
\textbf{Proposed partial credit.}\par\smallskip
\begin{tabular}{@{}l r@{}}\toprule
\CreditDescription{\textbf{Assessed step or criterion}} & \textbf{Points}\\\midrule
\CreditDescription{Variable, units and support} & 0.125\\
\CreditDescription{The denominator and exhaustive composition counts} & 0.5\\
\CreditDescription{The bonus, cap and aggregation} & 0.375\\
\CreditDescription{The grouped PMF table} & 0.25\\
\CreditDescription{The labeled, correctly spaced graph} & 0.25\\
\midrule
\CreditDescription{\textbf{Total}} & \textbf{1.5}\\
\bottomrule\end{tabular}
\end{minipage}\par\smallskip

\par\vspace{3mm}\secInicio[][sec:q6]

\clearpage
\textbf{Question 6 --- continued}\quad\hyperref[sec:q6]{Original description}

\Needspace{8\baselineskip}
\subsection*{Task 2 --- question and worked solution (proposed allocation) \score{3}}
\label{sol:q6t2}

\noindent\textbf{\boldmath Construct the complete \HLCumulativeProbability{CDF} table, including both tails, and draw a labeled step graph. Use open and filled endpoints correctly and preserve actual monetary spacing. Verify the defining \HLCumulativeProbability{CDF} properties and that its jumps equal the \GTenHLProbability{PMF} masses.}\par\nopagebreak

For $\HLCumulativeProbability{F_X(x)}=P(\GTenHLValue{X}\leq \GTenHLValue{x})$, add masses in increasing payout order. The successive cumulative numerators are $\HLCumulativeProbability{4},\allowbreak\HLCumulativeProbability{22},\allowbreak\HLCumulativeProbability{34},\allowbreak\HLCumulativeProbability{46},\allowbreak\HLCumulativeProbability{47},\allowbreak\HLCumulativeProbability{75},\allowbreak\HLCumulativeProbability{84}$, all over $\HLCumulativeProbability{84}$.

\begin{center}
\begin{tabular}{cc}\toprule
Cutoff interval (\GTenHLValue{USD}) & $\HLCumulativeProbability{F_X(x)}$\\\midrule
$\GTenHLValue{x}<\GTenHLValue{0}$ & $\HLCumulativeProbability{0}$\\
$\GTenHLValue{0}\leq \GTenHLValue{x}<\GTenHLValue{5}$ & $\HLCumulativeProbability{1/21}$\\
$\GTenHLValue{5}\leq \GTenHLValue{x}<\GTenHLValue{10}$ & $\HLCumulativeProbability{11/42}$\\
$\GTenHLValue{10}\leq \GTenHLValue{x}<\GTenHLValue{12}$ & $\HLCumulativeProbability{17/42}$\\
$\GTenHLValue{12}\leq \GTenHLValue{x}<\GTenHLValue{15}$ & $\HLCumulativeProbability{23/42}$\\
$\GTenHLValue{15}\leq \GTenHLValue{x}<\GTenHLValue{24}$ & $\HLCumulativeProbability{47/84}$\\
$\GTenHLValue{24}\leq \GTenHLValue{x}<\GTenHLValue{25}$ & $\HLCumulativeProbability{25/28}$\\
$\GTenHLValue{x}\geq \GTenHLValue{25}$ & $\HLCumulativeProbability{1}$\\
\bottomrule\end{tabular}
\end{center}

\begin{center}\begin{tikzpicture}
\begin{axis}[keyaxis,width=.86\linewidth,height=7cm,xmin=-2,xmax=27,ymin=-.07,ymax=1.1,xtick={0,5,10,12,15,24,25},xticklabel style={rotate=45,anchor=north east},ytick={0,.25,.5,.75,1},ylabel={$\HLCumulativeProbability{F_X(x)}$},yticklabel={$\HLCumulativeProbability{\pgfmathprintnumber{\tick}}$}]

\addplot[keyCDFColor,thick,<-] coordinates {(-2,0)(0,0)};

\addplot[keyCDFColor,thick] coordinates {(0,4/84)(5,4/84)};

\addplot[keyCDFColor,thick] coordinates {(5,22/84)(10,22/84)};

\addplot[keyCDFColor,thick] coordinates {(10,34/84)(12,34/84)};

\addplot[keyCDFColor,thick] coordinates {(12,46/84)(15,46/84)};

\addplot[keyCDFColor,thick] coordinates {(15,47/84)(24,47/84)};

\addplot[keyCDFColor,thick] coordinates {(24,75/84)(25,75/84)};

\addplot[keyCDFColor,thick,->] coordinates {(25,1)(27,1)};

\addplot[keyCDFColor,only marks,mark=*,mark size=.8pt,mark options={fill=white}] coordinates {(0,0/84)(5,4/84)(10,22/84)(12,34/84)(15,46/84)(24,47/84)(25,75/84)};

\addplot[keyCDFColor,only marks,mark=*,mark size=.8pt] coordinates {(0,4/84)(5,22/84)(10,34/84)(12,46/84)(15,47/84)(24,75/84)(25,84/84)};

\end{axis}\end{tikzpicture}\end{center}

\textbf{CDF checks.}
\begin{itemize}
\item Domain: the table covers every real cutoff.
\item Range: all tabulated heights lie in $[\HLCumulativeProbability{0},\HLCumulativeProbability{1}]$.
\item Monotonicity: the tabulated heights never decrease.
\item Right-continuity: each horizontal segment includes its left endpoint and excludes its right endpoint; the filled point gives the value at a jump.
\item Lower-tail limit: the left arrow extends to $-\infty$, where the limit is 0.
\item Upper-tail limit: the right arrow extends to $+\infty$, where the limit is 1.
\item Jump--mass agreement: in \GTenHLValue{support} order, subtracting successive cumulative numerators (starting at zero) gives $\GTenHLProbability{4},\allowbreak\GTenHLProbability{18},\allowbreak\GTenHLProbability{12},\allowbreak\GTenHLProbability{12},\allowbreak\GTenHLProbability{1},\allowbreak\GTenHLProbability{28},\allowbreak\GTenHLProbability{9}$. Hence each jump is its \GTenHLProbability{PMF} mass over $\GTenHLProbability{84}$: $\HLCumulativeProbability{F_X(a)}-\HLCumulativeProbability{F_X(a^-)}=\GTenHLProbability{p_X(a)}$. At unsupported cutoffs both the jump and the mass are zero.
\end{itemize}

\par\smallskip
\noindent\begin{minipage}{\linewidth}\centering
\textbf{Proposed partial credit.}\par\smallskip
\begin{tabular}{@{}l r@{}}\toprule
\CreditDescription{\textbf{Assessed step or criterion}} & \textbf{Points}\\\midrule
\CreditDescription{Cumulative sums} & 0.75\\
\CreditDescription{The complete interval table and tails} & 0.75\\
\CreditDescription{The labeled, correctly spaced step graph and endpoints} & 0.75\\
\CreditDescription{The CDF and jump checks} & 0.75\\
\midrule
\CreditDescription{\textbf{Total}} & \textbf{3}\\
\bottomrule\end{tabular}
\end{minipage}\par\smallskip

\par\vspace{3mm}\secInicio[][sec:q6]

\clearpage
\textbf{Question 6 --- continued}\quad\hyperref[sec:q6]{Original description}

\Needspace{8\baselineskip}
\subsection*{Task 3 (c) --- question and worked solution (proposed allocation) \score{1.5}}
\label{sol:q6t3}

\noindent\textbf{\boldmath Calculate and interpret (i) the probability that the \GTenHLValue{final payout} is less than \GTenHLValue{12 USD} and (ii) the probability that it is at least \GTenHLValue{24 USD}.}\par\nopagebreak

\textbf{Allocation:} (i) 0.75 points; (ii) 0.75 points.

\textbf{(i) Less than \GTenHLValue{12 USD}.} Include only the final-payout \GTenHLValue{support} values $\GTenHLValue{0},\GTenHLValue{5},\GTenHLValue{10}$; exclude the mass at \GTenHLValue{12 USD} because the inequality is strict:
\[
P(\GTenHLValue{X}<\GTenHLValue{12})=\GTenHLProbability{p_X(0)}+\GTenHLProbability{p_X(5)}+\GTenHLProbability{p_X(10)}
=\frac{\GTenHLProbability{4}+\GTenHLProbability{18}+\GTenHLProbability{12}}{\GTenHLProbability{84}}=\GTenHLProbability{\frac{17}{42}}\approx\GTenHLProbability{0.404762}.
\]
The \GTenHLValue{final payout} for a uniformly selected three-ticket subset is below \GTenHLValue{12 USD} with probability $\GTenHLProbability{17/42}$, approximately $\GTenHLProbability{40.48\%}$.

\textbf{(ii) At least \GTenHLValue{24 USD}.} Include both final payouts $\GTenHLValue{24}$ and $\GTenHLValue{25}$, including the boundary mass at \GTenHLValue{24 USD} and the capped mass at \GTenHLValue{25 USD}:
\[
P(\GTenHLValue{X}\geq\GTenHLValue{24})=\GTenHLProbability{p_X(24)}+\GTenHLProbability{p_X(25)}
=\frac{\GTenHLProbability{28}+\GTenHLProbability{9}}{\GTenHLProbability{84}}=\GTenHLProbability{\frac{37}{84}}\approx\GTenHLProbability{0.440476}.
\]
The \GTenHLValue{final payout} for a uniformly selected three-ticket subset is \GTenHLValue{24 USD} or more with probability $\GTenHLProbability{37/84}$, approximately $\GTenHLProbability{44.05\%}$, after the bonus and cap.

\par\smallskip
\noindent\begin{minipage}{\linewidth}\centering
\textbf{Proposed partial credit.}\par\smallskip
\begin{tabular}{@{}l r@{}}\toprule
\CreditDescription{\textbf{Assessed step or criterion}} & \textbf{Points}\\\midrule
\CreditDescription{(i): selecting exactly the included PMF masses with the correct boundary} & 0.375\\
\CreditDescription{(i): summing them correctly} & 0.1875\\
\CreditDescription{(i): the contextual interpretation} & 0.1875\\
\CreditDescription{(ii): selecting exactly the included PMF masses with the correct boundary} & 0.375\\
\CreditDescription{(ii): summing them correctly} & 0.1875\\
\CreditDescription{(ii): the contextual interpretation} & 0.1875\\
\midrule
\CreditDescription{\textbf{Total}} & \textbf{1.5}\\
\bottomrule\end{tabular}
\end{minipage}\par\smallskip

Each part totals 0.75 points.

\par\vspace{3mm}\secInicio[][sec:q6]

\textbf{Scoring reconciliation.} Question 6 gives $1.5+3+1.5=6$ points for C2. C2 gives $2+2+4+6=14$ points; its independent passing threshold is 4 points. C1 gives $3+3=6$ points; its independent passing threshold is also 4 points. The overall available total is $6+14=20$ points, with no transfer between criteria.

\end{document}
